% Part: normal-modal-logic
% Chapter: axioms-systems
% Section: duals

\documentclass[../../../include/open-logic-section]{subfiles}

\begin{document}

\olfileid{nml}{prf}{dua}

\olsection{દ્વૈત \usetoken{P}{formula}}

\begin{defn}\ollabel{def:duals}
  !!{formula}s \Ax{T}, \Ax{B}, \Ax{4} અને \Ax{5} દરેકનું
  \emph{દ્વૈત} છે; તેને નીચેની જેમ અધોલિખિત હીરાથી દર્શાવીએ છીએ:
  \begin{align}
    \tag{\Ax{T_\Diamond}} p & \lif \Diamond p\\
    \tag{\Ax{B_\Diamond}} \Diamond\Box p & \lif p\\
    \tag{\Ax{4_\Diamond}} \Diamond\Diamond p & \lif \Diamond p\\
    \tag{\Ax{5_\Diamond}} \Diamond\Box p & \lif \Box p
    \end{align}
\end{defn}

ઉપરના દરેક દ્વૈત !!{formula}ને અનુરૂપ !!{formula}માંથી મેળવવા પહેલાં
$p$ના સ્થાને $\lnot p$ મૂકો, પછી પ્રતિપક્ષી વિધાન લો, પછી
$\lnot\Box\lnot$ના સ્થાને $\Diamond$ અને $\lnot\Diamond\lnot$ના
સ્થાને~$\Box$ મૂકો. આ અર્થમાં \Ax{D}, એટલે કે
$\Box!A \lif \Diamond!A$, પોતાનું જ દ્વૈત છે.

\begin{prop}\ollabel{prop:dualsys}
  \olref[nml][prf][dua]{def:duals}માંના દરેક !!{formula}~$!A$ માટે:
  $\Log{K}!A = \Log{K}!A_{\Diamond}$.
\end{prop}
\begin{proof}
  અભ્યાસપ્રશ્ન.
\end{proof}
\begin{prob}
  \olref[nml][prf][dua]{prop:dualsys} સાબિત કરો.
\end{prob}

\end{document}
